\section{Safeguarding the Right to Data Protection as a Fundamental Right}

Good administration provides the basis for efficient, fair, and rational administrative decision making, as well as guaranteeing individual procedural rights.\autocite[213]{hofmannDemkovaPrinciplesProcJustice}
Compliance with the standards set by good administration upholds the integrity of the public administrative system, as well as an individual's understanding of the reasons and input of a decision.\autocite[213]{hofmannDemkovaPrinciplesProcJustice}

This study examined the reforms introduced by the Procedural Regulation insofar as they relate to the functioning of the \gls{gdpr}'s cooperation mechanism, a network of peer authorities that are tasked with a uniform application of the law protecting the fundamental right to data protection.
So far, this thesis focused on the role of the actors involved in cross-border enforcement of the \gls{gdpr}, and their ability to push each other to comply with the standards of good administration through the tools of administrative cooperation provided in data protection law.
This chapter will consider whether good administration, as a principle, is an effective tool for safeguarding the right to data protection as a fundamental right, and what consequences stem from a failure to comply with its standards in the context of that fundamental right.

\subsection{Process Rights and Effective Protection for Individuals}

The institution of procedural rights is deeply connected to the respect of human dignity.\autocite[75--76]{galliganProceduralFairness1997}
Fair procedures bring fair treatment, and fair treatment is a necessity if we asssume that individuals are entitled to dignity and respect, especially in a community that is based on the rule of law.\autocites[75--76]{galliganProceduralFairness1997}[338--339]{nehlGoodAdminProceduralRight2009}
Procedural guarantees may not always be enough to ensure effective protection for the individual and collective, but function as a framework to assess the powers of the public authority, and seek redress when those powers are abused.\autocite[76--77]{galliganProceduralFairness1997}

In this context, good administration provides the judiciary with a flexible framework to calibrate the intensity of judicial oversight.\autocite[102]{hofmannDutyCareEU2020}
Judicial oversight, however, is limited by the separation of powers.\autocite[4]{mendesDiscretionCarePublic2016}
Review is not a substitute for administrative discretion, and the judiciary cannot command the choices made within the adminsitration's discretionary power.\autocite[4]{mendesDiscretionCarePublic2016}

\subsection{Is Good Administration an Effective Tool for Safeguarding the Right to Data Protection?}
While good administration has served as a benchmark for the assessment of the amended cross-border enforcement system provided by the Procedural Regulation, it is not a self-executing tool for ensuring the protection of the fundamental right to data protection.
The principles examined in this study -- namely the duty of the administration to deliver careful and reasoned decisions -- if violated, may form grounds for annullment of the administrative decision.\autocites[29]{TUM}[78]{Sytraval}
Annulment forms an individual remedy, and is not a substitute for the procedures meant to ensure the delicate balance between the protection of fundamental rights and the free movement of data established by the \gls{gdpr}.

\subsection{The Consequences of a Failure to Enforce the Right to Data Protection}

The right to data protection provides us with a right to informational self-determination, especially in light of the consequences the processing of personal data may bring in relation fundamental freedoms, such as the right to private life.\autocites[68]{ÖsterreichischerRundfunk}[7]{CFR}[8]{echr}
Data protection is not an absolute right,\autocite[recital 4]{GDPR} but processing must always be carried out in accordance with the law,\autocite[8(2)]{CFR} to ensure that the balance between the protection of fundamental rights and the interests requiring the free movement of data is fair.\autocite[42]{SchremsI}
The guarantee of independence of the supervisory authority is intended to ensure the effectiveness and reliability of the system of data protection, and supervision is therefore an essential component of the protection of individuals against infringements of their fundamental right to data protection.\autocite[41]{SchremsI}

Striking a balance between the protection of fundamental rights and the free movement of data is no easy feat.
Advancements in technology -- and in particular more recently, the use of neural networks -- have enabled the processing of personal data at a scale unimaginable even a decade ago.\autocite[recital 6]{GDPR}
For this reason, the \gls{gdpr} created a system of principles and rules meant to be technologically neutral, to ensure the law cannot be circumvented.\autocite[recital 15]{GDPR}
Many new pieces of legislation refer to the \gls{gdpr}, or are without prejudice to it, requiring a careful examination of the interactions with more specific legislation, and how these intesections will shape the emerging digital acquis.\autocite[215]{zanfirGDPREDPS}

A disjointed system of oversight creates `hawks and doves', leading to forum shopping and an evasion of the law.\autocites[800]{gentileDeficientDesign2022}[point 124]{fbBelgiumAgBobek}
When violations of the \gls{gdpr} are not investigated to their fullest, and the checks and balances of the integrated system are prevented from functioning, the protection of the fundamental right to data protection is undermined.

This may have real consequences for individuals and society at large, especially since technology and data have taken over virtually all aspects of our lives.\autocite{stevensSphereTransgressionsReflecting2024}
Violations of the right to data protection may lead to losses of autonomy, a degradation of democracy, and violations of other fundamental rights, such as freedom of expression or non-discrimination.\autocite[19, 52--53, 348]{fraHandbookDataprotection2018}
Data protection authorities have a duty to ensure that the law is enforced, deterring and punishing actors who violate the fundamental right to data protection.